\section{Threats to validity, and what this run could not measure}
\label{sec:threats}

\noindent\textbf{The binaries are tuned to this machine.}
Compilation uses the native architecture flags, so absolute cycle counts belong
to the microarchitecture they were measured on. Portable binaries would
understate what a tuned data plane achieves, which is the wrong error for this
study to make. The crossover $s^{*}$ is a ratio and transfers; $a$ and $b$
separately do not.

\noindent\textbf{The host is shared and virtualised.}
The measurement host is classified \envHostClass, and every caveat that
classification implies travels with the dataset. Where kernel isolation on the
data-plane cores is absent, another tenant's work appears as variance rather
than as bias. That is what the interleaved sweep order and the anchor check are
for: drift that would otherwise land on one factor level is spread across all of
them, and a machine that changed under the measurement voids the sweep instead
of reporting it.

\noindent\textbf{The path crosses each byte twice.}
The data plane decrypts and re-encrypts. That is a real forwarding topology and
it makes $b$ large enough to separate cleanly from $a$, but it is not the only
topology. A one-pass endpoint would have roughly half this $b$ and therefore
roughly twice this $s^{*}$. The per-pass figures are reported alongside so that
a reader can make that adjustment rather than guess at it.

\noindent\textbf{The bound on kernel bypass is conditional.}
A submission poller that shares a core with the data plane is a competitor, not
a bypass. Where that condition holds, the difference between the polled and
unpolled paths bounds nothing, and Section~\ref{sec:iopath} reports the
condition rather than the bound.

\noindent\textbf{Bare metal is not available.}
The virtualisation tiers are measured above a hypervisor guest, so they give the
incremental cost of each isolation layer and not the cost of virtualisation
itself. On this run two of the four tiers did not build, and the two that did
ordered themselves the wrong way round, so no per-tier penalty is claimed.

\noindent\textbf{Three experiments did not run here.}
The cross-node scaling fragment, the latency curve and the GPU offload are all
absent, and each is recorded as unavailable with its reason rather than omitted.
The consequences are specific. Without the cross-node fragment, the claim that
per-packet cost is unchanged by putting ranks on separate hosts is untested.
Without the latency curve, the capacity in Section~\ref{sec:latency} implies a
knee that was not observed, and nothing in this report characterises behaviour
above it. Without the GPU arm, the constants of Equation~\ref{eq:gpu} are
unmeasured, and the argument that the bus decides the offload rests on the shape
of the model rather than on this machine.

\noindent\textbf{The admission model has no data to fit.}
The measurements were intended to support a control question: given the offered
rate, the payload, the cipher and the network conditions, can the p99 latency be
predicted well enough to decide whether to admit more traffic? The comparison
that would matter is not whether a model fits, since a model fitted to a smooth
curve always fits. It is whether a learned predictor beats a static threshold
that admits traffic up to a fixed fraction of measured capacity, with splits
taken by network condition rather than shuffled, so that neighbouring offered
rates from one condition cannot leak across the boundary. That comparison needs
the latency points the previous paragraph says are missing, so it is not
reported.

\noindent\textbf{One machine.}
Every number here comes from one host. The merge step refuses to combine
fragments from different machines or builds, precisely because that would
otherwise be invisible, but the consequence is that generality across hardware
is argued from the structure of the model and not demonstrated.
